\textbf{Phần 1. Trắc nghiệm lựa chọn}
\setcounterpageref{ex}{0}

\begin{ex} %[3D1V01]
Số tự nhiên $X$ có ba chữ số. Biết chữ số hàng trăm là số chẵn lớn nhất có một chữ số, chữ số hàng chục kém chữ số hàng trăm $2$ đơn vị, chữ số hàng đơn vị bằng một nửa chữ số hàng chục. Cách đọc của số $X$ là
\choice
{Tám trăm sáu mươi}
{Tám trăm ba mươi sáu}
{\True Tám trăm sáu mươi ba}
{Sáu trăm tám mươi ba}
\loigiai{
Chữ số chẵn lớn nhất có một chữ số là $8$, do đó chữ số hàng trăm là $8$.\\
Chữ số hàng chục là: $8 - 2 = 6$.\\
Chữ số hàng đơn vị là: $6 : 2 = 3$.\\
Vậy số $X$ được viết là $863$ và đọc là: \textit{Tám trăm sáu mươi ba}.\\
\textbf{Đáp án của câu hỏi là C}
}
\end{ex}

\begin{ex} %[3D1V01]
Trong số $949$, giá trị của chữ số $9$ ở hàng trăm gấp giá trị của chữ số $9$ ở hàng đơn vị số lần là
\choice
{\True $100$ lần}
{$10$ lần}
{$90$ lần}
{$1000$ lần}
\loigiai{
Chữ số $9$ ở hàng trăm có giá trị là $900$.\\
Chữ số $9$ ở hàng đơn vị có giá trị là $9$.\\
Giá trị chữ số $9$ ở hàng trăm gấp giá trị chữ số $9$ ở hàng đơn vị số lần là:
$$900 : 9 = 100\text{ (lần)}$$
\textbf{Đáp án của câu hỏi là A}
}
\end{ex}

\begin{ex} %[3D1V01]
Số liền sau của số chẵn bé nhất có ba chữ số khác nhau là
\choice
{$100$}
{$102$}
{$101$}
{\True $103$}
\loigiai{
Để số chẵn có 3 chữ số khác nhau là bé nhất:\\
- Chữ số hàng trăm bé nhất khác 0 là $1$.\\
- Chữ số hàng chục bé nhất khác 1 là $0$.\\
- Chữ số hàng đơn vị là số chẵn bé nhất khác 0 là $2$.\\
Số chẵn bé nhất có 3 chữ số khác nhau là $102$.\\
Số liền sau của $102$ là: $102 + 1 = 103$.\\
\textbf{Đáp án của câu hỏi là D}
}
\end{ex}

\begin{ex} %[3D1V01]
Một số gồm $5$ trăm, $14$ chục và $6$ đơn vị. Số đó được viết là
\choice
{$5146$}
{\True $646$}
{$546$}
{$656$}
\loigiai{
Ta có:\\
- $5$ trăm $= 500$\\
- $14$ chục $= 140$\\
- $6$ đơn vị $= 6$\\
Số đó là: $500 + 140 + 6 = 646$.\\
\textbf{Đáp án của câu hỏi là B}
}
\end{ex}

\begin{ex} %[3D1V01]
Cho biểu thức $A = 600 + 30 + 8$ và $B = 600 + 50 + 2$. Số tròn chục $x$ thỏa mãn $A < x < B$ là
\choice
{$630$}
{$650$}
{\True $640$}
{$660$}
\loigiai{
Ta có: $A = 638$ và $B = 652$.\\
Điều kiện: $638 < x < 652$.\\
Vì $x$ là số tròn chục nên $x = 640$.\\
\textbf{Đáp án của câu hỏi là C}
}
\end{ex}

\begin{ex} %[3D1V01]
Số lẻ lớn nhất có ba chữ số khác nhau mà chữ số hàng trăm gấp $3$ lần chữ số hàng đơn vị là
\choice
{\True $983$}
{$981$}
{$973$}
{$389$}
\loigiai{
Vì số cần tìm là số lẻ nên chữ số hàng đơn vị phải là số lẻ: $1, 3, 5, 7, 9$.\\
Mặt khác, chữ số hàng trăm gấp 3 lần chữ số hàng đơn vị nên chữ số hàng đơn vị chỉ có thể là 1 (hàng trăm là 3) hoặc 3 (hàng trăm là 9).\\
Để số là lớn nhất, ta chọn chữ số hàng trăm là 9 và chữ số hàng đơn vị là 3.\\
Chữ số hàng chục lớn nhất khác 9 và 3 là 8.\\
Vậy số cần tìm là $983$.\\
\textbf{Đáp án của câu hỏi là A}
}
\end{ex}

\begin{ex} %[3D1V01]
Cho các giá trị sau:\\
$M = 700 + 80 + 5$;\\
$N$ là số liền trước của số $790$;\\
$P$ là số gồm $8$ trăm và $5$ đơn vị;\\
$Q = 700 + 95$.\\
Sắp xếp các số $M, N, P, Q$ theo thứ tự từ bé đến lớn là
\choice
{$M < Q < N < P$}
{$N < M < Q < P$}
{$M < N < P < Q$}
{\True $M < N < Q < P$}
\loigiai{
Tính giá trị từng số:\\
- $M = 785$\\
- $N = 790 - 1 = 789$\\
- $P = 805$\\
- $Q = 795$\\
Vì $785 < 789 < 795 < 805$ nên thứ tự từ bé đến lớn là: $M < N < Q < P$.\\
\textbf{Đáp án của câu hỏi là D}
}
\end{ex}

\begin{ex} %[3D1V01]
Một số có ba chữ số có chữ số hàng trăm là $5$. Nếu xóa bỏ chữ số $5$ này thì số đó giảm đi bao nhiêu đơn vị?
\choice
{$50$ đơn vị}
{\True $500$ đơn vị}
{$5$ đơn vị}
{$450$ đơn vị}
\loigiai{
Gọi số ban đầu có dạng $\overline{5ab} = 500 + \overline{ab}$.\\
Khi xóa chữ số 5 ở hàng trăm, ta được số mới là $\overline{ab}$.\\
Số đó giảm đi số đơn vị là:
$$(500 + \overline{ab}) - \overline{ab} = 500\text{ (đơn vị)}$$
\textbf{Đáp án của câu hỏi là B}
}
\end{ex}

\begin{ex} %[3D1V01]
Từ ba chữ số $0; 4; 7$, lập được bao nhiêu số chẵn có ba chữ số khác nhau?
\choice
{$4$ số}
{$2$ số}
{\True $3$ số}
{$1$ số}
\loigiai{
Một số là số chẵn khi chữ số hàng đơn vị là 0 hoặc 4.\\
- \textbf{Trường hợp 1:} Chữ số hàng đơn vị là 0. Các số lập được là: $470; 740$ (có 2 số).\\
- \textbf{Trường hợp 2:} Chữ số hàng đơn vị là 4. Chữ số hàng trăm không thể là 0 nên chỉ có thể là 7. Số lập được là: $704$ (có 1 số).\\
Tổng cộng có: $2 + 1 = 3$ số thỏa mãn.\\
\textbf{Đáp án của câu hỏi là C}
}
\end{ex}

\begin{ex} %[3D1V01]
Hiệu giữa số lớn nhất có 3 chữ số khác nhau và số bé nhất có 3 chữ số khác nhau được lập từ ba chữ số $3; 0; 8$ là
\choice
{\True $522$}
{$530$}
{$502$}
{$492$}
\loigiai{
- Số lớn nhất có 3 chữ số khác nhau lập từ $3; 0; 8$ là: $830$.\\
- Số bé nhất có 3 chữ số khác nhau lập từ $3; 0; 8$ là: $308$.\\
Hiệu hai số là: $830 - 308 = 522$.\\
\textbf{Đáp án của câu hỏi là A}
}
\end{ex}

\begin{ex} %[3D1V01]
Cho dãy số cách đều: $120; 125; 130; 135; \dots$ Số hạng thứ 10 của dãy số là
\choice
{$160$}
{$170$}
{$175$}
{\True $165$}
\loigiai{
Dãy số có khoảng cách giữa hai số liền nhau là: $125 - 120 = 5$ đơn vị.\\
Số hạng thứ 10 của dãy số là:
$$120 + (10 - 1) \times 5 = 120 + 45 = 165$$
\textbf{Đáp án của câu hỏi là D}
}
\end{ex}

\begin{ex} %[3D1V01]
Tìm số có ba chữ số, biết chữ số hàng đơn vị là $2$, chữ số hàng chục gấp $3$ lần chữ số hàng đơn vị và chữ số hàng trăm kém chữ số hàng chục $2$ đơn vị. Số đó là
\choice
{$642$}
{\True $462$}
{$426$}
{$264$}
\loigiai{
Chữ số hàng đơn vị là $2$.\\
Chữ số hàng chục là: $2 \times 3 = 6$.\\
Chữ số hàng trăm là: $6 - 2 = 4$.\\
Vậy số cần tìm là $462$.\\
\textbf{Đáp án của câu hỏi là B}
}
\end{ex}

\textbf{Phần 2. Trắc nghiệm đúng sai}
\setcounterpageref{ex}{0}

\begin{ex} %[3D1V01]
Xét tính đúng sai của các nhận định sau về các số trong phạm vi 1000:
\choiceTF
{\True Hiệu giữa số lớn nhất có ba chữ số và số bé nhất có ba chữ số là $899$.}
{Số liền trước của một số có ba chữ số luôn luôn là một số có ba chữ số.}
{\True Trong số $444$, giá trị của chữ số $4$ ở hàng trăm gấp $100$ lần giá trị của chữ số $4$ ở hàng đơn vị.}
{\True Có tất cả $90$ số tròn chục có ba chữ số.}
\loigiai{
\begin{itemize}
    \item Ý a) Đúng: Số lớn nhất có 3 chữ số là 999, số bé nhất có 3 chữ số là 100. Hiệu là $999 - 100 = 899$.
    \item Ý b) Sai: Số liền trước của số 100 là 99 (là số có 2 chữ số).
    \item Ý c) Đúng: Chữ số 4 hàng trăm có giá trị 400, chữ số 4 hàng đơn vị có giá trị 4. Ta có $400 : 4 = 100$.
    \item Ý d) Đúng: Các số tròn chục có 3 chữ số từ 100 đến 990. Số lượng số là: $(990 - 100) : 10 + 1 = 90$ số.
\end{itemize}
\begin{itemize}
    \item a) \textbf{Đúng}.
    \item b) \textbf{Sai}.
    \item c) \textbf{Đúng}.
    \item d) \textbf{Đúng}.
\end{itemize}
}
\end{ex}

\begin{ex} %[3D1V01]
Cho tập hợp ba chữ số: $2; 5; 0$. Xét tính đúng sai của các khẳng định sau:
\choiceTF
{\True Từ ba chữ số đã cho lập được tất cả 4 số có ba chữ số khác nhau.}
{Trong các số có ba chữ số khác nhau lập được, chỉ có đúng 2 số chia hết cho 5.}
{\True Số lẻ duy nhất có ba chữ số khác nhau lập được từ ba chữ số trên là $205$.}
{\True Tổng của số lớn nhất và số bé nhất có ba chữ số khác nhau lập được bằng $725$.}
\loigiai{
\begin{itemize}
    \item Ý a) Đúng: Các số lập được là $205; 250; 502; 520$ (4 số).
    \item Ý b) Sai: Các số có tận cùng là 0 hoặc 5 gồm $205; 250; 520$ (có 3 số chia hết cho 5).
    \item Ý c) Đúng: Chỉ có số $205$ có tận cùng là 5 (số lẻ).
    \item Ý d) Đúng: Số lớn nhất là 520, số bé nhất là 205. Tổng là: $520 + 205 = 725$.
\end{itemize}
\begin{itemize}
    \item a) \textbf{Đúng}.
    \item b) \textbf{Sai}.
    \item c) \textbf{Đúng}.
    \item d) \textbf{Đúng}.
\end{itemize}
}
\end{ex}

\textbf{Phần 3. Trả lời ngắn}
\setcounterpageref{ex}{0}

\begin{ex} %[3D1V01]
Tìm số tự nhiên có ba chữ số, biết tổng ba chữ số bằng $4$ và chữ số hàng trăm lớn hơn chữ số hàng chục, chữ số hàng chục lớn hơn chữ số hàng đơn vị.
\shortans{310}
\loigiai{
\textbf{Đáp số:} $310$\\
Gọi số cần tìm là $\overline{abc}$ với $a > b > c \ge 0$ và $a + b + c = 4$.\\
Vì $a > b > c$ nên $c$ bé nhất là 0. Khi $c = 0$, ta có $a + b = 4$ với $a > b > 0$. Ta chọn được $a = 3, b = 1$.\\
Vậy số cần tìm là $310$.
}
\end{ex}

\begin{ex} %[3D1V01]
Dãy số chẵn liên tiếp $102; 104; 106; \dots; 180$ có tất cả bao nhiêu số hạng?
\shortans{40}
\loigiai{
\textbf{Đáp số:} $40$\\
Số số hạng của dãy là: $(180 - 102) : 2 + 1 = 78 : 2 + 1 = 40$ (số hạng).
}
\end{ex}

\begin{ex} %[3D1V01]
Để đánh số trang của một cuốn sách từ trang $1$ đến trang $100$, người ta phải dùng tất cả bao nhiêu chữ số?
\shortans{192}
\loigiai{
\textbf{Đáp số:} $192$\\
- Từ trang 1 đến trang 9 có 9 trang có 1 chữ số, cần: $9 \times 1 = 9$ (chữ số).\\
- Từ trang 10 đến trang 99 có: $(99 - 10) + 1 = 90$ trang có 2 chữ số, cần: $90 \times 2 = 180$ (chữ số).\\
- Trang 100 là trang có 3 chữ số, cần: $3$ (chữ số).\\
Tổng số chữ số cần dùng là: $9 + 180 + 3 = 192$ (chữ số).
}
\end{ex}

\begin{ex} %[3D1V01]
Tìm số lớn nhất có ba chữ số khác nhau mà tổng ba chữ số bằng $14$ và chữ số hàng trăm là $8$.
\shortans{860}
\loigiai{
\textbf{Đáp số:} $860$\\
Gọi số cần tìm có dạng $\overline{8ab}$ với $a \ne b \ne 8$.\\
Tổng hai chữ số hàng chục và hàng đơn vị là: $14 - 8 = 6$.\\
Để số là lớn nhất thì chữ số hàng chục $a$ phải lớn nhất có thể ($a \le 6$). Chọn $a = 6$, khi đó chữ số hàng đơn vị là $b = 6 - 6 = 0$.\\
Vậy số lớn nhất cần tìm là $860$.
}
\end{ex}

\begin{ex} %[3D1V01]
Cho ba chữ số $1; 5; 9$. Tính tổng của số lớn nhất và số bé nhất có ba chữ số khác nhau được lập từ ba chữ số đã cho.
\shortans{1110}
\loigiai{
\textbf{Đáp số:} $1110$\\
- Số lớn nhất có ba chữ số khác nhau lập từ $1; 5; 9$ là: $951$.\\
- Số bé nhất có ba chữ số khác nhau lập từ $1; 5; 9$ là: $159$.\\
Tổng hai số là: $951 + 159 = 1110$.
}
\end{ex}

\begin{ex} %[3D1V01]
Có bao nhiêu số có ba chữ số mà chữ số hàng trăm bằng chữ số hàng đơn vị?
\shortans{90}
\loigiai{
\textbf{Đáp số:} $90$\\
Số có dạng $\overline{aba}$ ($a \ne 0$):\\
- Chữ số hàng trăm $a$ có 9 cách chọn ($a \in \{1; 2; \dots; 9\}$).\\
- Chữ số hàng chục $b$ có 10 cách chọn ($b \in \{0; 1; 2; \dots; 9\}$).\\
- Chữ số hàng đơn vị phải bằng $a$ nên có đúng 1 cách chọn.\\
Số các số thỏa mãn là: $9 \times 10 \times 1 = 90$ (số).
}
\end{ex}
